\section{Methodology}\label{sec:method}


\subsection{Heuristics, key ideas and jump-augmented SDE}

In contrast to Brownian motion, non-Gaussian L\'{e}vy processes (e.g., $\alpha$-stable L\'{e}vy processes) are characterized by discontinuous trajectories with non-local jump displacements. These macroscopic spatial excursions significantly enhance transition probabilities across high energy barriers, effectively mitigating the metastability issue. Leveraging the L\'{e}vy score function framework \cite{huang2025probability,Huang2024LevyScore}, it is feasible to incorporate non-local jump processes into the underlying SDE \eqref{eqn:sde} while strictly preserving the target invariant Boltzmann distribution $p_\infty$. Specifically, by employing the L\'{e}vy score function, the evolution of the probability density $p_t$ can be reformulated via the following integro-differential L\'{e}vy--Fokker--Planck equation (LFPE):
\begin{equation}\label{eqn:LFPE}
    \begin{aligned}
        \partial_t p_t(x)  =& \nabla\cdot\big(\nabla V p_t\big)(x) + \frac{\sigma^2}{2} \triangle p_t(x)\\
        =& \nabla\cdot\left[\left(\nabla V(x) + \int_0^1\int_{\mathbb{R}^d-\{0\}} \frac{r p_t(x - \theta r)}{p_t(x)}\nu(\d r)\d\theta\right)p_t(x)\right] + \frac{\sigma^2}{2} \triangle p_t(x) \\
        &\quad + \int_{\mathbb{R}^d-\{0\}}\big(p_t(x-r) -p_t(x)\big)\nu(\d r),
    \end{aligned}
\end{equation}
where $\nu$ represents an appropriate L\'{e}vy measure. Here, the expression
\begin{equation}\label{Levyscore}
    S_\mathrm{L}(t,x):= -\int_0^1\int_{\mathbb{R}^d-\{0\}} \frac{r p_t(x - \theta r)}{p_t(x)}\nu(\d r)\d\theta,
\end{equation}
defines the L\'{e}vy score function \cite{Huang2024LevyScore}, a pivotal concept in the analysis of probability flow and entropy production  for jump-diffusion processes \cite{huang2025entropy}. Treating the L\'{e}vy score function $S_\mathrm{L}(t,x)$ as an a priori known term, the LFPE structure implies that the probability flow $p_t$ corresponds to the following associated SDE (cf. \cite[Theorem 2.3]{huang2025probability}):
\begin{equation}\label{eqn:sdeLevy}
    \d Y_t=\left(-\nabla V(Y_{t-}) - \int_0^1\int_{\mathbb{R}^d-\{0\}} \frac{r p_t(Y_{t-} - \theta r)}{p_t(Y_{t-})}\nu(\d r)\d\theta \right)\d t+ \sigma \d B_t + \d L_t,\quad Y_0\sim p_0,
\end{equation}
where $L_t$ is a L\'{e}vy process independent of both the Brownian motion $B_t$ and the initial distribution $p_0$, driven by a L\'{e}vy measure $\nu$ governed by the score function $S_\mathrm{L}$ in \eqref{Levyscore}.

Although the diffusion process $X_t$ in \eqref{eqn:sde} and the jump-diffusion process $Y_t$ in \eqref{eqn:sdeLevy} exhibit distinct pathwise behaviors, they inherently share the identical probability flow. Consequently, utilizing the time-dependent jump-diffusion SDE \eqref{eqn:sdeLevy} to sample the invariant measure offers no theoretical acceleration compared to the diffusion SDE \eqref{eqn:sde}. This limitation stems from the density-dependent L\'{e}vy score function $S_\mathrm{L}(t,x)$ effectively counteracting the macroscopic jumps of the L\'{e}vy process $L_t$ within the probability space. However, as the system relaxes toward its stationary state, the L\'{e}vy score converges to a time-independent asymptotic form:
\begin{equation} \label{eqn:stationaryLevyScore}
\begin{aligned}
    S^\mathrm{s}_\mathrm{L}(x):=& -\int_0^1\int_{\mathbb{R}^d-\{0\}} \frac{r p_\infty(x - \theta r)}{p_\infty(x)}\nu(\d r)\d\theta \\
    =& -\int_0^1\int_{\mathbb{R}^d-\{0\}} r\exp\left\{-\frac{2V(x- \theta r)}{\sigma^2} + \frac{2V(x)}{\sigma^2}\right\} \nu(\d r)\d\theta.
    \end{aligned}
\end{equation}
This observation motivates the proposal of the following autonomous jump-diffusion SDE for accelerated sampling purposes:
\begin{equation}\label{eqn:sdeLevy2}
(\mbox{jump-argmented SDE})\quad\quad    \d Z_t=\left(-\nabla V(Z_{t-}) + S^\mathrm{s}_\mathrm{L}\big( Z_{t-}\big) \right)\d t+ \sigma \d B_t + \d L_t.
\end{equation}

Prior to discussing the numerical sampling strategy, we establish the theoretical properties of the proposed dynamics \eqref{eqn:sdeLevy2}. To facilitate the analysis, we require additional structural properties for the potential function $V$, as detailed in the following assumption.

\begin{assumption}\label{assumption:U}
     The potential $V$ is coercive, satisfying $V(x)\to\infty$ as $|x|\to\infty$. Furthermore, we impose a monotonicity condition on the gradient magnitude: there exists a constant $R>0$ such that
     \[
     |\nabla V(x_1)|\geq|\nabla V(x_2)| \quad \text{whenever} \quad |x_1|\geq|x_2|>R.
     \]
\end{assumption}

\begin{remark}
    Assumption \ref{assumption:U} is analogous to conditions employed in the study of $\Gamma$-convergence \cite{lu2017gaussian}. In the present context, this assumption is crucial for analyzing the well-posedness of the jump-diffusion SDE \eqref{eqn:sdeLevy2}, as it guarantees that the potential $V$ possesses sufficient geometric regularity and growth properties at infinity.
\end{remark}

To address technical considerations regarding potentially unbounded drifts, we analyze the jump-diffusion SDE \eqref{eqn:sdeLevy2} through a sequence of truncated approximations. For any fixed $M>0$, let $h:\mathbb{R}\to\mathbb{R}$ be a smooth scalar cutoff function defined by
\begin{equation}
    h(s):=\begin{cases}
        1,\quad s\leq M-1,\\
        0,\quad s\geq M.
    \end{cases}
\end{equation}
We introduce an auxiliary truncated potential $V^M$ associated with $V$ as follows:
\begin{equation}\label{VM}
    V^M(x):= h\big(U(x)\big)U(x) + \left[1-h\big (U(x)\big )\right]|x|^2,\quad x\in\mathbb{R}^d.
\end{equation}
Consider the corresponding SDE governed by this truncated potential:
\begin{equation}\label{eqn:sdeM}
    \d X^M_t= -\nabla V^M\big(X^M_t\big)\d t + \sigma \d B_t.
\end{equation}
It is straightforward to verify that the modified drift term satisfies the coercivity and monotonicity properties outlined in Assumption \ref{assum:SDE}. The associated invariant Boltzmann distribution for this truncated system reads
\begin{equation}\label{p^M_infty}
    p^M_\infty(x)=\big(Z^M\big)^{-1}\exp\left\{-\frac{2V^M(x)}{\sigma^2}\right\}.
\end{equation}
Consequently, the global well-posedness and the existence of a unique invariant measure for the SDE \eqref{eqn:sdeM} are guaranteed. Analogous to the derivation in the preceding section, we define the truncated jump-diffusion SDE corresponding to \eqref{eqn:sdeLevy2}, which recovers the original dynamics in the limit as $M\to\infty$:
\begin{equation}\label{eqn:sdeLevy2M}
    \d Z^M_t=\left(-\nabla V^M\big(Z^M_{t-}\big) + S^{\mathrm{s}M}_\mathrm{L}\big( Z^M_{t-}\big) \right)\d t+ \sigma \d B_t + \d L_t,
\end{equation}
where
\begin{equation}
\begin{aligned}
    S^{\mathrm{s}M}_\mathrm{L}(x):=& -\int_0^1\int_{\mathbb{R}^d-\{0\}} r\exp\left\{-\frac{2V^M(x- \theta r)}{\sigma^2} + \frac{2V^M(x)}{\sigma^2}\right\} \nu(\d r)\d\theta.
    \end{aligned}
\end{equation}

\subsection{Global well-posedness and ergodicity of jump-argmented SDE \texorpdfstring{\eqref{eqn:sdeLevy2}}{(2.1)}}

We proceed under the following additional structural assumptions regarding the L\'{e}vy measure $\nu$.

\begin{assumption}\label{assumption}
    The L\'{e}vy measure $\nu$ possesses a finite second moment, i.e., $\int_{\mathbb{R}^d-\{0\}}|r|^2\nu(\d r)<\infty$. 
\end{assumption}

\begin{theorem}\label{theo:well-posedness}
    Under Assumptions \ref{assum:SDE}, \ref{assumption:U}, and \ref{assumption}, the jump-diffusion SDE \eqref{eqn:sdeLevy2} is globally well-posed.
\end{theorem}

\begin{proof}
Observe that for $x\in\mathbb{R}^d$ satisfying $V(x)>M$, the drift term involving the modified L\'{e}vy score $S^{\mathrm{s}M}_\mathrm{L}(x)$ satisfies the following estimate:
\begin{equation}
    \begin{aligned}
        & -2\int_0^1\int_{\mathbb{R}^d-\{0\}} r\cdot x\exp\left\{-\frac{2V^M\big(x - \theta r\big)}{\sigma^2} + \frac{2V^M\big(x \big)}{\sigma^2}\right\} \nu(\d r)\d\theta\\
        =& -2\int_0^1\int_{\mathbb{R}^d-\{0\}} r\cdot x\exp\left\{-\frac{2 \big|x - \theta r\big|^2}{\sigma^2} + \frac{2|x|^2}{\sigma^2} \right\} \nu(\d r)\d\theta\\
        =& -2\int_0^1\int_{\mathbb{R}^d-\{0\}} r\cdot x\exp\left\{ \frac{4 x \cdot \theta r - 2|\theta r|^2}{\sigma^2}  \right\} \nu(\d r)\d\theta\\
         \leq & 8|x|^2\int_0^1\int_{\mathbb{R}^d-\{0\}} \frac{  |\theta r|^2}{\sigma^2}\exp\left\{-  \frac{2|\theta r|^2}{\sigma^2} \right\} \nu(\d r)\d\theta.
    \end{aligned}
\end{equation}
Because the potential $V$ is coercive (Assumption \ref{assumption:U}), the sublevel set $\{x\in\mathbb{R}^d\mid V(x)\leq M\}$ is compact. Consequently, the continuous function $x\cdot S^{\mathrm{s}M}_\mathrm{L}(x)$ is bounded on this set, yielding
\begin{equation}
    \sup_{\{x\in\mathbb{R}^d\mid V(x)\leq M\}} \left| x\cdot S^{\mathrm{s}M}_\mathrm{L}(x)\right|<\infty.
\end{equation}
Let us define the constant $K_1'$ as
\begin{equation}
    K_1':= K_1 + \sup_{\{x\in\mathbb{R}^d\mid V(x)\leq M\}} \left| x\cdot S^{\mathrm{s}M}_\mathrm{L}(x)\right| + 8 \int_0^1\int_{\mathbb{R}^d-\{0\}} \frac{  |\theta r|^2}{\sigma^2}\exp\left\{-  \frac{2|\theta r|^2}{\sigma^2} \right\} \nu(\d r)\d\theta<\infty,
\end{equation}
where $K_1$ is the constant introduced in Assumption \ref{assum:SDE}. Then, for all $x\in\mathbb{R}^d$, we obtain the global bound
\begin{equation}
    2x\cdot \left(-\nabla V^M(x) + S^{\mathrm{s}M}_\mathrm{L}(x)\right)\leq K'_1(|x|^2 + 1).
\end{equation}
Thus, by invoking \cite[Theorem 2.2]{xi2019jump}, we conclude that the solution process of the truncated system \eqref{eqn:sdeLevy2M} is non-explosive. Furthermore, it is straightforward to verify that the modified drift $-\nabla V^M + S^{\mathrm{s}M}_\mathrm{L}$ is locally Lipschitz continuous. The local Lipschitz continuity of the drift term guarantees the uniqueness of the solution to \eqref{eqn:sdeLevy2M}. Taking the limit as $M\to\infty$, this sequence of well-posed local approximations systematically extends to a globally well-posed solution for the original jump-diffusion SDE \eqref{eqn:sdeLevy2}, which completes the proof.
\end{proof}

To further guarantee that the jump-diffusion SDE \eqref{eqn:sdeLevy2} possesses well-defined transition densities and a unique invariant distribution, we introduce the following regularity assumption.

\begin{assumption}\label{ass:main}
    \hfill
    \begin{enumerate}
        \item The potential $V$ is infinitely differentiable, i.e., $V \in C^\infty(\mathbb{R}^d,\mathbb{R})$.
        \item There exists a constant $\delta > 0$ such that  
        \[
            \limsup_{\epsilon \to 0} \frac{\left(\int_{|r| \leq \epsilon} |r|^3 \nu(\d r)\right)^{1/3}}{\left(\int_{|r| \leq \epsilon} |r|^2 \nu(\d r)\right)^{(1 + \delta)/2}} < \infty.
        \]\label{subassum:nu}
    \end{enumerate}
\end{assumption}

\begin{remark}
    Assumption \ref{ass:main} is introduced to mathematically ensure the existence of smooth transition densities for the jump-diffusion SDE \eqref{eqn:sdeLevy2}. Alternative, simpler conditions can also be employed to achieve this. For instance, consider the scenario where the L\'{e}vy process is a compound Poisson process governed by the L\'{e}vy measure $\lambda_0\nu_J$, where $\lambda_0$ acts as the intensity parameter and $\nu_J$ characterizes the distribution of the jump size. In this case, it is sufficient for $\nu_J$ to possess bounded support and a smooth density $\nu_J\in C^\infty(\mathbb{R}^d,\mathbb{R})$; for further details and proofs, we refer the reader to \cite[Theorem 2-29]{bichteler1987malliavin}.
\end{remark}

\begin{theorem}\label{theo:p_infty}
    Under Assumptions \ref{assum:SDE}, \ref{assumption:U}, \ref{assumption}, and \ref{ass:main}, the jump-diffusion SDE \eqref{eqn:sdeLevy2} admits a unique invariant measure $p_\infty$, which precisely corresponds to the Boltzmann distribution of the SDE \eqref{eqn:sde}, in the sense that for every compact set $\mathcal{K}\subset\mathbb{R}^d$,
    \begin{equation}
        \lim_{M\to\infty}\sup_{x\in\mathcal{K}}|p^M_\infty(x)-p_\infty(x)|=0.
    \end{equation}
\end{theorem}

\begin{proof}
    By construction, the truncated auxiliary potential $V^M$ defined in \eqref{VM} is infinitely differentiable. Furthermore, although the drift vector field $-\nabla V^M$ is not necessarily globally bounded, all of its higher-order derivatives remain bounded. Consequently, under Assumptions \ref{assum:SDE}, \ref{assumption}, and \ref{ass:main}, the solution process $Z^M_t$ to the truncated jump-diffusion SDE \eqref{eqn:sdeLevy2M} admits unique $C^\infty$ transition probability densities, a regularity property rigorously established in \cite[Theorem 3.4]{bally2024upper}. This inherent regularity ensures that the stationary distribution $p^M_\infty$ defined in \eqref{p^M_infty} acts as the unique invariant measure for the truncated process $Z^M_t$. Specifically, this holds because $p^M_\infty$ solves the stationary LFPE associated with \eqref{eqn:sdeLevy2M}:
    \begin{equation}
        \begin{aligned}
            0 =& \nabla\cdot\left[\left(\nabla V^M(x) + \int_0^1\int_{\mathbb{R}^d-\{0\}}r\exp\left\{-\frac{2V^M(x- \theta r)}{\sigma^2} + \frac{2V^M(x)}{\sigma^2}\right\} \nu(\d r)\d\theta \right)p^M_\infty(x)\right]  \\
            &\quad + \frac{\sigma^2}{2} \triangle p^M_\infty(x) + \int_{\mathbb{R}^d-\{0\}}\big(p^M_\infty(x-r) -p^M_\infty(x)\big)\nu(\d r)\\
            =&\nabla\cdot \left[\left(\nabla V^M(x)    + \frac{\sigma^2}{2} \nabla\log p^M_\infty(x)\right)p^M_\infty(x)\right].
        \end{aligned}
    \end{equation}
    Taking the limit as $M\to\infty$ yields the desired uniform convergence on compact sets to the target distribution $p_\infty$, thereby completing the proof.
\end{proof}



\subsection{Spectral analysis for jump-augmented SDE \eqref{eqn:sdeLevy2}}
In this section, we proceed to analyze the underlying theoretical mechanisms driving the accelerated sampling rate of the proposed methodology.
%=======================================================
\subsubsection{Closed-form Dirichlet form of jump-augmented SDE \eqref{eqn:sdeLevy2}}
The following assumption is imposed \emph{only} to justify the explicit energy identity below under the non-truncated convention \eqref{eq:generator_levy}.

\begin{assumption}[Finite activity]\label{ass:finite_activity_dirichlet}
The L\'evy measure satisfies $\nu(\mathbb{R}^d\setminus\{0\})<\infty$.
\end{assumption}


\begin{remark}[On truncation conventions]\label{rem:truncation_convention}
Assumption \ref{ass:finite_activity_dirichlet} (compound Poisson/finite activity) together with Assumption \ref{assumption}
ensures that the non-truncated jump generator in \eqref{eq:generator_levy} is well-defined on $C_c^\infty$
and justifies repeated use of Tonelli/Fubini in the proof below.
If one instead works with the standard L\'evy--It\^{o} truncation
$\int(\varphi(x+r)-\varphi(x)-\nabla\varphi(x)\cdot r\,\mathbf{1}_{\{|r|\le 1\}})\nu(\mathrm{d}r)$,
then the forward equation and the energy identity acquire the corresponding compensator terms.
\end{remark}

%--------------------------------------------------------
%\paragraph{Main identity (energy decomposition).}
Under Assumptions \ref{assumption} and \ref{ass:finite_activity_dirichlet}, it can be shown that
the L\'evy-score drift of jump-diffusion SDE \eqref{eqn:sdeLevy2} cancels the non-symmetric part induced by the pure-jump generator at the level of the quadratic form.
Consequently, the energy admits a closed-form decomposition into a diffusion part and a jump increment part.

\begin{proposition}[Closed-form Dirichlet energy]\label{thm:closed_form_energy}
Assume Assumptions \ref{assumption} and \ref{ass:finite_activity_dirichlet}. 
Let $\mathcal{A}$ be given by \eqref{eq:generator_levy} with
$b(x)=-\nabla V(x)+S_{\mathrm L}^{\mathrm s}(x)$, where $S_{\mathrm L}^{\mathrm s}$ is defined by \eqref{eqn:stationaryLevyScore} and $\mathcal{E}_{\mathrm{levy}}(f,g):=-\langle f,\mathcal{A}g\rangle_{p_\infty}$.
Then for any $f\in C_c^\infty(\mathbb{R}^d)$,
\begin{equation}\label{eq:closed_form_energy}
\mathcal{E}_{\mathrm{levy}}(f,f)
=
\frac{\sigma^2}{2}\int_{\mathbb{R}^d}|\nabla f(x)|^2\,p_\infty(x)\,\mathrm{d}x
+\frac12\iint_{\mathbb{R}^d\times(\mathbb{R}^d\setminus\{0\})}\big(f(x+r)-f(x)\big)^2\,p_\infty(x)\,\nu(\mathrm{d}r)\,\mathrm{d}x.
\end{equation}
In particular, $\mathcal{E}_{\mathrm{levy}}(f,f)\ge 0$.
\end{proposition}

\begin{remark}\label{rem:dissipation_equals_energy}
For any $f\in C_c^\infty(\R^d)$, the quadratic form in Proposition~\ref{thm:closed_form_energy}
satisfies $\mathcal E_{\mathcal A}(f,f)=-\mathrm{Re}\langle f,\mathcal A f\rangle_{p_\infty}
=\mathcal E_{\mathrm{levy}}(f,f)\ge 0$.
\end{remark}

\begin{proof}
For any $f\in C_c^\infty(\mathbb{R}^d)$, by definition, \ref{def:dirichletlevy} we know that 
\begin{align*}
\mathcal{E}_{\mathrm{levy}}(f,f)
% &=-\langle f,\mathcal{A}f\rangle_{p_\infty}\\
&= -\int_{\mathbb{R}^d} f(x)\,b(x)\cdot\nabla f(x)\,p_\infty(x)\,\mathrm{d}x
-\frac{\sigma^2}{2}\int_{\mathbb{R}^d} f(x)\,\Delta f(x)\,p_\infty(x)\,\mathrm{d}x\\
&\quad\;-\int_{\mathbb{R}^d} f(x)\,(\mathcal{J}f)(x)\,p_\infty(x)\,\mathrm{d}x.
\end{align*}

\emph{Step 1: diffusion contribution.}
Using $p_\infty(x)\propto \exp\{-2V(x)/\sigma^2\}$ and integration by parts (valid since $f$ is compactly supported), the potential drift
$-\nabla V$ cancels with the $\Delta$-term to yield
\[
-\int f(-\nabla V)\cdot\nabla f\,p_\infty\,\mathrm{d}x
-\frac{\sigma^2}{2}\int f\,\Delta f\,p_\infty\,\mathrm{d}x
=
\frac{\sigma^2}{2}\int |\nabla f|^2\,p_\infty\,\mathrm{d}x.
\]

\emph{Step 2: jump term and score cancellation.}
By expansion of $\mathcal{J}$,
\begin{equation}\label{eq:jump_expansion_in_proof}
-\int f(x)\,(\mathcal{J}f)(x)\,p_\infty(x)\,\mathrm{d}x
=\iint \big(f(x)^2-f(x)f(x+r)\big)\,p_\infty(x)\,\nu(\mathrm{d}r)\,\mathrm{d}x.
\end{equation}
Next, by the definition  of $S_{\mathrm{L}}^{\mathrm{s}}$ \eqref{eqn:stationaryLevyScore} and Tonelli's theorem (Assumptions \ref{ass:finite_activity_dirichlet} and \ref{assumption}),
\begin{align*}
-\int f(x)\,S_{\mathrm{L}}^{\mathrm{s}}(x)\cdot \nabla f(x)\,p_\infty(x)\,\mathrm{d}x
&=\int_{\mathbb{R}^d}\int_0^1\int_{\mathbb{R}^d\setminus\{0\}}
f(x)\, r\cdot \nabla f(x)\,p_\infty(x-\theta r)\,\nu(\mathrm{d}r)\,\mathrm{d}\theta\,\mathrm{d}x.
\end{align*}
For fixed $r\neq 0$ and $\theta\in[0,1]$, set $y=x-\theta r$ to obtain
\[
\int_{\mathbb{R}^d} f(x)\, r\cdot \nabla f(x)\,p_\infty(x-\theta r)\,\mathrm{d}x
=\int_{\mathbb{R}^d} f(y+\theta r)\, r\cdot \nabla f(y+\theta r)\,p_\infty(y)\,\mathrm{d}y.
\]
Integrating over $\theta\in[0,1]$ and using
\[
\frac{\mathrm{d}}{\mathrm{d}\theta}\,\frac12 f(y+\theta r)^2
= f(y+\theta r)\, r\cdot \nabla f(y+\theta r),
\]
we obtain
\[
\int_0^1 f(y+\theta r)\, r\cdot \nabla f(y+\theta r)\,\mathrm{d}\theta
=\frac12\big(f(y+r)^2-f(y)^2\big).
\]
Therefore,
\begin{equation}\label{eq:score_term_in_proof}
-\int f\,S_{\mathrm{L}}^{\mathrm{s}}\cdot \nabla f\,p_\infty\,\mathrm{d}x
=\frac12\iint \big(f(x+r)^2-f(x)^2\big)\,p_\infty(x)\,\nu(\mathrm{d}r)\,\mathrm{d}x.
\end{equation}
Summing \eqref{eq:jump_expansion_in_proof} and \eqref{eq:score_term_in_proof} gives
\[
-\int f\,S_{\mathrm{L}}^{\mathrm{s}}\cdot \nabla f\,p_\infty\,\mathrm{d}x
-\int f\,(\mathcal{J}f)\,p_\infty\,\mathrm{d}x
=\frac12\iint \big(f(x+r)-f(x)\big)^2\,p_\infty(x)\,\nu(\mathrm{d}r)\,\mathrm{d}x.
\]
Combining this with Step 1 yields \eqref{eq:closed_form_energy}.
\end{proof}

%--------------------------------------------------------
For completeness, we define the symmetric part of the full bilinear form as follows.
\begin{definition}[Symmetric bilinear form]
The symmetric part of the full bilinear form is given by
\[
\mathcal{E}^{\mathrm{sym}}_{\mathrm{levy}}(f,g)
:=\frac12\Big(\mathcal{E}_{\mathrm{levy}}(f,g)+\mathcal{E}_{\mathrm{levy}}(g,f)\Big),
\qquad f,g\in C_c^\infty(\mathbb{R}^d).
\]
\end{definition}
Then Proposition \ref{thm:closed_form_energy} implies
\[
\mathcal{E}^{\mathrm{sym}}_{\mathrm{levy}}(f,f)
=\mathcal{E}_{\mathrm{levy}}(f,f)
=\frac{\sigma^2}{2}\int_{\mathbb{R}^d}|\nabla f|^2\,p_\infty\,\mathrm{d}x
+\frac12\iint (f(x+r)-f(x))^2\,p_\infty(x)\,\nu(\mathrm{d}r)\,\mathrm{d}x.
\]




%========================================================
\subsubsection{\(\lambda(\mathcal A)\ge \lambda(\mathcal L)\)}

\[
\mathcal E_{\mathrm{levy}}(f,f)
=
\mathcal E_{\mathrm{diff}}(f,f)+\mathcal E_{\mathrm{jump}}(f,f),
\]
where
\[
\mathcal E_{\mathrm{jump}}(f,f)
:=
\frac12\int_{\mathbb R^d}\int_{\mathbb R^d\setminus\{0\}}
\big(f(x+r)-f(x)\big)^2\,p_\infty(x)\,\nu(dr)\,dx
\ge 0.
\]
for any \(h\neq 0\), we have
\[
\frac{\mathcal E_{\mathrm{levy}}(h,h)}{\|h\|^2}
\ge
\frac{\mathcal E_{\mathrm{diff}}(h,h)}{\|h\|^2}.
\]
So
\[
\lambda(\mathcal A)\ge \lambda(\mathcal L).
\]
%========================================================
% ============================================================
% Partition of State Space (two-well, barrier-free via jumps)
% Notation consistent with the original setup:
%   - invariant density: p_\infty
%   - diffusion energy:  \mathcal E_{\mathrm{diff}}
%   - jump energy:       \mathcal E_{\mathrm{jump}}
%   - global jump:       \nu_{\mathrm{glob}} = \lambda_0 \nu_J
% ============================================================
\subsection{2-well}
\subsubsection{Partition of State Space}\label{sec:partition}

\paragraph{Basins of attraction and geometric partition.}
Assume that the energy landscape $V:\mathbb R^d\to\mathbb R$ is of double-well type and admits exactly two distinct local minima (well bottoms) $m_1,m_2\in\mathbb R^d$.
Consider the deterministic gradient flow
\[
\dot x_t \;=\; -\nabla V(x_t),
\qquad x_0 = x\in\mathbb R^d,
\]
and denote by $x_t(x)$ the (maximal) solution starting from $x$.
The (forward) basin of attraction of $m_i$ is defined by
\[
\widetilde\Omega_i
\;:=\;
\Bigl\{x\in\mathbb R^d:\ x_t(x)\ \text{exists for all }t\ge 0
\ \text{and}\ \lim_{t\to\infty}x_t(x)=m_i\Bigr\},
\qquad i\in\{1,2\}.
\]
We assume that there are no other recurrent attractors so that, up to a negligible set of initial conditions,
each point flows to one of the two minima.
More precisely, there exists a Borel set $\mathcal N\subset\mathbb R^d$ with
\[
\mathcal N\subset\partial\widetilde\Omega_1\cap\partial\widetilde\Omega_2,
\qquad
\mathcal N\ \text{has Lebesgue and $p_\infty$-measure zero},
\]
such that
\[
\mathbb R^d
\;=\;
\widetilde\Omega_1\cup\widetilde\Omega_2\cup\mathcal N,
\qquad
\widetilde\Omega_1\cap\widetilde\Omega_2=\varnothing.
\]
For definiteness we remove the null set and work with the measurable partition
\[
\Omega_i \;:=\; \widetilde\Omega_i\setminus\mathcal N,
\qquad i\in\{1,2\},
\]
so that
\[
\Omega_1\cup\Omega_2=\mathbb R^d,
\qquad
\Omega_1\cap\Omega_2=\varnothing.
\]



\paragraph{Invariant measure and well-weights.}
Let $p_\infty$ denote the invariant density on $\mathbb R^d$ and define the well-weights
\[
w_i \;:=\; \int_{\Omega_i} p_\infty(x)\,\mathrm dx,
\qquad i\in\{1,2\}.
\]
Clearly,
\[
w_1+w_2=1.
\]

\paragraph{Conditional invariant measures on each region.}
Define the conditional densities (and the associated probability measures) by
\[
p_\infty^{(i)}(x)
\;:=\;
\frac{1}{w_i}\,p_\infty(x)\,\mathbf 1_{\Omega_i}(x),
\qquad
p_\infty^{(i)}(\mathrm dx) := p_\infty^{(i)}(x)\,\mathrm dx.
\]
For any integrable test function $f:\mathbb R^d\to\mathbb R$, define the conditional means
\[
\bar f_{\Omega_i}
\;:=\;
\int_{\Omega_i} f(x)\,p_\infty^{(i)}(\mathrm dx).
\]
Define the conditional variances
\[
\mathrm{Var}_{p_\infty^{(i)}}(f)
\;:=\;
\int_{\Omega_i}\bigl(f(x)-\bar f_{\Omega_i}\bigr)^2\,p_\infty^{(i)}(\mathrm dx).
\]

\paragraph{Exact variance decomposition (two-state coarse graining).}
With the above partition, the global variance admits the exact decomposition
\begin{equation}\label{eq:var-decomp-Omega}
\mathrm{Var}_{p_\infty}(f)
\;=\;
\sum_{i=1}^2 w_i\,\mathrm{Var}_{p_\infty^{(i)}}(f)
\;+\;
w_1w_2\bigl(\bar f_{\Omega_1}-\bar f_{\Omega_2}\bigr)^2.
\end{equation}
This is the (exact) ``within-well + between-well'' decomposition:
the first term measures fluctuations inside each region $\Omega_i$,
while the second term measures the mismatch of the two regional averages.
\subsubsection{Local Cores Inside Each Well}\label{sec:cores}

\paragraph{Local minima and core sets.}
Let $m_1,m_2\in\mathbb R^d$ denote the two local minima (well bottoms) and define their distance
\[
D \;:=\; |m_2-m_1|.
\]
Fix a core radius $r_{\mathrm c}>0$ and define the core sets
\[
\mathcal C_i \;:=\; B(m_i,r_{\mathrm c}),
\qquad i\in\{1,2\}.
\]
Assume the cores lie in their respective regions:
\[
\mathcal C_1\subset \Omega_1,
\qquad
\mathcal C_2\subset \Omega_2,
\]
and (for convenience in the shell argument) assume also $D>4r_{\mathrm c}$ so that $\mathcal C_1\cap\mathcal C_2=\varnothing$.

\paragraph{Core masses and core-conditional measures.}
Define the core masses
\[
c_i \;:=\; \int_{\mathcal C_i} p_\infty(x)\,\mathrm dx,
\qquad i\in\{1,2\},
\]
and define the core-conditional densities
\[
p_{\infty,\mathcal C}^{(i)}(x)
\;:=\;
\frac{1}{c_i}\,p_\infty(x)\,\mathbf 1_{\mathcal C_i}(x),
\qquad
p_{\infty,\mathcal C}^{(i)}(\mathrm dx) := p_{\infty,\mathcal C}^{(i)}(x)\,\mathrm dx.
\]
Define the core means
\[
\bar f_{\mathcal C_i}
\;:=\;
\int_{\mathcal C_i} f(x)\,p_{\infty,\mathcal C}^{(i)}(\mathrm dx).
\]
Finally define the core fractions inside each region:
\[
\eta_i \;:=\; p_\infty^{(i)}(\mathcal C_i) \;=\; \frac{c_i}{w_i}\in(0,1].
\]

% ============================================================
% Dirichlet forms and spectral gap
% ============================================================

\subsubsection{Dirichlet Forms and Spectral Gap}\label{sec:dirichlet-gap}

\paragraph{Diffusion and jump Dirichlet forms.}
Let $\sigma>0$ denote the diffusion scale. Define
\[
\mathcal E_{\mathrm{diff}}(f,f)
\;:=\;
\frac{\sigma^2}{2}\int_{\mathbb R^d} |\nabla f(x)|^2\,p_\infty(x)\,\mathrm dx,
\]
and for a jump measure $\nu$ define
\[
\mathcal E_{\mathrm{jump}}(f,f)
\;:=\;
\frac12\int_{\mathbb R^d} p_\infty(x)
\int_{\mathbb R^d}\bigl(f(x+r)-f(x)\bigr)^2\,\nu(\mathrm dr)\,\mathrm dx.
\]
The total Dirichlet form is
\[
\mathcal E_{\mathrm{levy}}(f,f) \;:=\; \mathcal E_{\mathrm{diff}}(f,f)+\mathcal E_{\mathrm{jump}}(f,f).
\]

\paragraph{Spectral gap via Poincar\'e inequality.}
Define the spectral gap $\lambda(\mathcal G)$ by the best constant in
\[
\mathrm{Var}_{p_\infty}(f)
\;\le\;
\frac{1}{\lambda(\mathcal G)}\,\mathcal E_{\mathrm{levy}}(f,f).
\]

% ============================================================
% Within-well control: \lambda_{\mathrm{in}}
% ============================================================


\subsubsection{Within-Well Mixing Rate \texorpdfstring{$\lambda_{\mathrm{in}}$}{lambda_in}}
\label{sec:within-well}



\paragraph{Restricted Dirichlet forms inside each well.}
Recall that the diffusion Dirichlet form is
\[
\mathcal E_{\mathrm{diff}}(f,f)
\;=\;
\frac{\sigma^2}{2}\int_{\mathbb R^d} |\nabla f(x)|^2\,p_\infty(x)\,\mathrm dx.
\]
For each basin \(\Omega_i\) we define the within-well diffusion Dirichlet form
by restricting both the domain and the reference measure:
\[
\mathcal E_{\mathrm{diff}}^{(i)}(f,f)
\;:=\;
\frac{\sigma^2}{2}\int_{\Omega_i} |\nabla f(x)|^2\,p_\infty^{(i)}(\mathrm dx),
\qquad
i\in\{1,2\}.
\]
The corresponding within-well spectral gap is
\[
\lambda_{\mathrm{in}}^{(i)}
\;:=\;
\inf\Bigl\{
\frac{\mathcal E_{\mathrm{diff}}^{(i)}(f,f)}{\mathrm{Var}_{p_\infty^{(i)}}(f)}
:\ f\in C_c^\infty(\Omega_i),\ \mathrm{Var}_{p_\infty^{(i)}}(f)>0
\Bigr\},
\]
that is, \(\lambda_{\mathrm{in}}^{(i)}\) is the best constant in the Poincar\'e
inequality
\begin{equation}\label{eq:local-Poincare-well}
\mathrm{Var}_{p_\infty^{(i)}}(f)
\;\le\;
\frac{1}{\lambda_{\mathrm{in}}^{(i)}}\,\mathcal E_{\mathrm{diff}}^{(i)}(f,f),
\qquad
\forall f\in C_c^\infty(\Omega_i).
\end{equation}
We then define the overall within-well mixing rate as
\[
\lambda_{\mathrm{in}}
\;:=\;
\min\bigl\{\lambda_{\mathrm{in}}^{(1)},\lambda_{\mathrm{in}}^{(2)}\bigr\}.
\]

\paragraph{Strong convexity on the core sets.}
Recall the core sets
\[
\mathcal C_i \;:=\; B(m_i,r_{\mathrm c}),
\qquad i\in\{1,2\},
\]
introduced in Section~\ref{sec:cores}.
We now assume that the potential is uniformly strongly convex on each core.

\begin{assumption}[Core strong convexity]\label{ass:core-strong-convex}
There exists a constant \(m_{\mathrm{in}}>0\) such that
\[
\nabla^2 V(x)
\;\succeq\;
m_{\mathrm{in}}\,\mathrm{Id}_d
\qquad
\text{for all }x\in\mathcal C_1\cup\mathcal C_2.
\]
\end{assumption}

Under Assumption~\ref{ass:core-strong-convex}, the invariant density
\(p_\infty\) is locally strongly log-concave on each \(\mathcal C_i\).
More precisely, for the diffusion
\[
\mathrm dX_t \;=\; -\nabla V(X_t)\,\mathrm dt + \sigma\,\mathrm dW_t,
\]
the invariant density is (up to normalization) \(p_\infty(x)\propto
\exp\bigl(-\tfrac{2}{\sigma^2}V(x)\bigr)\).
Thus on \(\mathcal C_i\) we have
\[
\nabla^2\!\Bigl(\tfrac{2}{\sigma^2}V(x)\Bigr)
\;\succeq\;
\frac{2m_{\mathrm{in}}}{\sigma^2}\,\mathrm{Id}_d.
\]
By the Brascamp--Lieb inequality, the core-conditional measures \(p_{\infty,\mathcal C}^{(i)}\)
satisfy a uniform local Poincar\'e inequality: for all smooth
\(g:\mathbb R^d\to\mathbb R\),
\begin{equation}\label{eq:core-Poincare}
\mathrm{Var}_{p_{\infty,\mathcal C}^{(i)}}(g)
\;\le\;
\kappa_{\mathrm{core}}
\int_{\mathcal C_i} |\nabla g(x)|^2\,p_{\infty,\mathcal C}^{(i)}(\mathrm dx),
\qquad
i\in\{1,2\},
\end{equation}
where we can take the explicit bound
\begin{equation}\label{eq:kappa-core-explicit}
\kappa_{\mathrm{core}}
\;\le\;
\frac{\sigma^2}{2m_{\mathrm{in}}}.
\end{equation}

\paragraph{Lyapunov drift outside the cores.}
We now impose a coercivity condition on \(V\) outside each core, along rays
emanating from the corresponding local minimum \(m_i\).
For \(x\in\mathbb R^d\) we write
\[
r_i(x) \;:=\; |x-m_i|,
\qquad i\in\{1,2\}.
\]

\begin{assumption}[Radial drift outside the core]\label{ass:drift-outside-core}
For each \(i\in\{1,2\}\) there exists a constant \(\alpha_i>0\) such that
\begin{equation}\label{eq:drift-condition-core}
\bigl\langle x-m_i,\nabla V(x)\bigr\rangle
\;\ge\;
\alpha_i\,r_i(x)
\qquad
\text{for all }x\in\Omega_i\ \text{with}\ r_i(x)\ge r_{\mathrm c}.
\end{equation}
\end{assumption}

We now construct an explicit Lyapunov function \(W_i\) on \(\Omega_i\).
Fix \(\gamma_i>0\) (to be chosen below) and set
\[
W_i(x)
\;:=\;
\begin{cases}
\exp\bigl(\gamma_i r_i(x)\bigr), & r_i(x)\ge r_{\mathrm c},\\[4pt]
\text{any $C^2$ extension with }W_i(x)\ge 1, & r_i(x)<r_{\mathrm c}.
\end{cases}
\]
By construction \(W_i\in C^2(\Omega_i)\) and \(W_i\ge 1\) everywhere.
On the exterior region \(\{x\in\Omega_i:\ r_i(x)\ge r_{\mathrm c}\}\), we have
\[
\nabla W_i(x)
=
\gamma_i\,\frac{x-m_i}{r_i(x)}\,W_i(x),
\qquad
\Delta W_i(x)
=
\Bigl(\gamma_i^2 + \gamma_i\frac{d-1}{r_i(x)}\Bigr)W_i(x).
\]
The diffusion generator reads
\[
L_{\mathrm{diff}} f(x)
\;:=\;
\frac{\sigma^2}{2}\Delta f(x) - \nabla V(x)\cdot\nabla f(x),
\]
and we compute, for \(r_i(x)\ge r_{\mathrm c}\),
\begin{align*}
L_{\mathrm{diff}} W_i(x)
&=
\frac{\sigma^2}{2}\Delta W_i(x) - \nabla V(x)\cdot\nabla W_i(x)\\
&=
\Biggl[
\frac{\sigma^2}{2}\Bigl(\gamma_i^2 + \gamma_i\frac{d-1}{r_i(x)}\Bigr)
- \gamma_i\,\frac{\langle x-m_i,\nabla V(x)\rangle}{r_i(x)}
\Biggr] W_i(x).
\end{align*}
Using Assumption~\ref{ass:drift-outside-core}, we obtain for all
\(x\in\Omega_i\) with \(r_i(x)\ge r_{\mathrm c}\),
\begin{align*}
L_{\mathrm{diff}} W_i(x)
&\le
\Biggl[
\frac{\sigma^2}{2}\Bigl(\gamma_i^2 + \gamma_i\frac{d-1}{r_i(x)}\Bigr)
- \gamma_i\alpha_i
\Biggr] W_i(x)\\
&\le
-\gamma_i
\Biggl[
\alpha_i
- \frac{\sigma^2}{2}\gamma_i
- \frac{\sigma^2}{2}\frac{d-1}{r_{\mathrm c}}
\Biggr] W_i(x).
\end{align*}
Assume that
\[
\alpha_i
\;>\;
\frac{\sigma^2}{2}\frac{d-1}{r_{\mathrm c}},
\]
and choose \(\gamma_i>0\) such that
\[
0<\gamma_i
\;<\;
\alpha_i
- \frac{\sigma^2}{2}\frac{d-1}{r_{\mathrm c}}.
\]
We then define the explicit drift constant
\begin{equation}\label{eq:theta-i-core-def}
\theta_i
\;:=\;
\gamma_i\Biggl(
\alpha_i
- \frac{\sigma^2}{2}\gamma_i
- \frac{\sigma^2}{2}\frac{d-1}{r_{\mathrm c}}
\Biggr)
\;>\;0
\end{equation}
and conclude that
\begin{equation}\label{eq:Lyap-outside-core}
L_{\mathrm{diff}} W_i(x)
\;\le\;
-\theta_i\,W_i(x),
\qquad
x\in\Omega_i,\ r_i(x)\ge r_{\mathrm c}.
\end{equation}

On the core region \(\mathcal C_i\) we define the finite constant
\begin{equation}\label{eq:b-i-core-def}
b_i
\;:=\;
\sup_{x\in\mathcal C_i}
\Bigl(
L_{\mathrm{diff}} W_i(x) + \theta_i W_i(x)
\Bigr)_+,
\qquad
(a)_+ := \max\{a,0\}.
\end{equation}
Since \(\mathcal C_i\) is compact and \(W_i\in C^2(\Omega_i)\), this supremum is finite.
Combining \eqref{eq:Lyap-outside-core} with \eqref{eq:b-i-core-def} yields the
global Lyapunov condition
\begin{equation}\label{eq:Lyap-full-core}
L_{\mathrm{diff}} W_i(x)
\;\le\;
-\theta_i\,W_i(x)
\;+\;
b_i\,\mathbf 1_{\mathcal C_i}(x),
\qquad
\forall x\in\Omega_i.
\end{equation}

\paragraph{Application of the BCG Lyapunov criterion.}
We now recall the relevant result of Bakry, Barthe, Cattiaux and Guillin
(BCG) that links Lyapunov conditions and global Poincar\'e inequalities.

\begin{theorem}[BCG Lyapunov criterion for Poincar\'e, informal statement]
\label{thm:BCG-informal}
Let \(\mu\) be a probability measure on \(\mathbb R^d\) which is reversible
for a diffusion generator \(L\) with carré du champ \(\Gamma\).
Assume there exist a measurable function \(W\ge 1\), constants
\(\theta>0, b<\infty\), and a measurable set \(K\subset\mathbb R^d\) such that
\[
LW \;\le\; -\theta W + b\,\mathbf 1_K
\quad\text{on }\mathbb R^d,
\]
and that the restriction of \(\mu\) to \(K\) satisfies a local Poincar\'e
inequality with constant \(C_P(K)\).
Then \(\mu\) satisfies a global Poincar\'e inequality
\[
\mathrm{Var}_\mu(f)
\;\le\;
\frac{1}{\theta}\bigl(1 + b\,C_P(K)\bigr)
\int \Gamma(f,f)\,\mathrm d\mu,
\qquad
\forall f\in C_c^\infty(\mathbb R^d).
\]
(See, for instance, Theorem~1.2 in Bakry--Barthe--Cattiaux--Guillin.)
\end{theorem}

We apply Theorem~\ref{thm:BCG-informal} to each conditional invariant
measure \(p_\infty^{(i)}\) with the following choices:
\[
\mu = p_\infty^{(i)},\qquad
L = L_{\mathrm{diff}},\qquad
\Gamma(f,f) = \frac{\sigma^2}{2}|\nabla f|^2,
\]
\[
W = W_i,\qquad
\theta = \theta_i,\qquad
b = b_i,\qquad
K = \mathcal C_i.
\]
The Lyapunov condition \eqref{eq:Lyap-full-core} provides the required drift
inequality, while the local Poincar\'e inequality on \(K=\mathcal C_i\) is
given by \eqref{eq:core-Poincare} with constant
\(C_P(K)=\kappa_{\mathrm{core}}\).

Therefore, Theorem~\ref{thm:BCG-informal} yields, for each
\(i\in\{1,2\}\) and every smooth \(f:\mathbb R^d\to\mathbb R\),
\[
\mathrm{Var}_{p_\infty^{(i)}}(f)
\;\le\;
\frac{1}{\theta_i}\bigl(1 + b_i\,\kappa_{\mathrm{core}}\bigr)
\int_{\Omega_i} \Gamma(f,f)\,p_\infty^{(i)}(\mathrm dx)
=
\frac{\sigma^2}{2}\frac{1 + b_i\,\kappa_{\mathrm{core}}}{\theta_i}
\int_{\Omega_i} |\nabla f(x)|^2\,p_\infty^{(i)}(\mathrm dx).
\]
Comparing with \eqref{eq:local-Poincare-well}, we obtain the explicit lower
bound on the within-well spectral gap
\begin{equation}\label{eq:lambda-in-i-lower-bound-core}
\lambda_{\mathrm{in}}^{(i)}
\;\ge\;
\frac{\sigma^2\,\theta_i}{2\bigl(1 + b_i\,\kappa_{\mathrm{core}}\bigr)},
\qquad
i\in\{1,2\},
\end{equation}
where \(\theta_i\) and \(b_i\) are given by
\eqref{eq:theta-i-core-def} and \eqref{eq:b-i-core-def}, and
\(\kappa_{\mathrm{core}}\) satisfies \eqref{eq:kappa-core-explicit}.

Finally, since
\[
\lambda_{\mathrm{in}}
\;:=\;
\min\bigl\{\lambda_{\mathrm{in}}^{(1)},\lambda_{\mathrm{in}}^{(2)}\bigr\},
\]
we arrive at the global within-well mixing bound
\begin{equation}\label{eq:lambda-in-global-core}
\lambda_{\mathrm{in}}
\;\ge\;
\min_{i\in\{1,2\}}
\frac{\sigma^2\,\theta_i}{2\bigl(1 + b_i\,\kappa_{\mathrm{core}}\bigr)},
\qquad
\kappa_{\mathrm{core}}
\;\le\;
\frac{\sigma^2}{2m_{\mathrm{in}}}.
\end{equation}
This shows that the within-well mixing rate is controlled explicitly by
the core strong convexity parameter \(m_{\mathrm{in}}\), the radial drift
constants \(\alpha_i\), the core radius \(r_{\mathrm c}\), and the scale
\(\sigma\) of the diffusion.

\subsubsection{Within-Well Mixing Rate \texorpdfstring{$\lambda_{\mathrm{in}}$}{lambda_in}}
\label{sec:within-well-v2}

\paragraph{Within-well Dirichlet forms and spectral gaps.}
Recall that the diffusion Dirichlet form is
\[
\mathcal E_{\mathrm{diff}}(f,f)
\;=\;
\frac{\sigma^2}{2}\int_{\mathbb R^d} |\nabla f(x)|^2\,p_\infty(x)\,\mathrm dx.
\]
For each basin \(\Omega_i\) we define the within-well diffusion Dirichlet form
by restricting both the domain and the reference measure:
\[
\mathcal E_{\mathrm{diff}}^{(i)}(f,f)
\;:=\;
\frac{\sigma^2}{2}\int_{\Omega_i} |\nabla f(x)|^2\,p_\infty^{(i)}(\mathrm dx),
\qquad
i\in\{1,2\},
\]
where \(p_\infty^{(i)}\) is the normalized restriction of \(p_\infty\) to \(\Omega_i\).
The associated within-well spectral gap is
\[
\lambda_{\mathrm{in}}^{(i)}
\;:=\;
\inf\Bigl\{
\frac{\mathcal E_{\mathrm{diff}}^{(i)}(f,f)}{\mathrm{Var}_{p_\infty^{(i)}}(f)}
:\ f\in C_c^\infty(\Omega_i),\ \mathrm{Var}_{p_\infty^{(i)}}(f)>0
\Bigr\},
\]
that is, \(\lambda_{\mathrm{in}}^{(i)}\) is the best constant in the Poincar\'e
inequality
\begin{equation}\label{eq:local-Poincare-well}
\mathrm{Var}_{p_\infty^{(i)}}(f)
\;\le\;
\frac{1}{\lambda_{\mathrm{in}}^{(i)}}\,\mathcal E_{\mathrm{diff}}^{(i)}(f,f),
\qquad
\forall f\in C_c^\infty(\Omega_i).
\end{equation}
We then define the overall within-well mixing rate as
\[
\lambda_{\mathrm{in}}
\;:=\;
\min\bigl\{\lambda_{\mathrm{in}}^{(1)},\lambda_{\mathrm{in}}^{(2)}\bigr\}.
\]

\paragraph{Radial drift outside the cores.}
For \(x\in\mathbb R^d\) and each \(i\in\{1,2\}\), write
\[
r_i(x) \;:=\; |x-m_i|,
\]
where \(m_i\) is the local minimum of \(V\) in \(\Omega_i\).
We impose the following coercive drift condition along rays emanating from \(m_i\).

\begin{assumption}[Radial drift outside the core]\label{ass:drift-outside-core}
There exist constants \(\alpha_i>0\) and \(r_{\mathrm c}>0\) such that, for each
\(i\in\{1,2\}\),
\begin{equation}\label{eq:drift-condition-core}
\bigl\langle x-m_i,\nabla V(x)\bigr\rangle
\;\ge\;
\alpha_i\,r_i(x)
\qquad
\text{for all }x\in\Omega_i\ \text{with}\ r_i(x)\ge r_{\mathrm c}.
\end{equation}
\end{assumption}

Note that we only assume \(\alpha_i>0\); \emph{no} additional relation between
\(\alpha_i\) and \(r_{\mathrm c}\) is required.

\paragraph{Enlarged balls and local Poincar\'e constants.}
Motivated by Assumption~\ref{ass:drift-outside-core}, we introduce a radius
\(R_i\) and the corresponding ball \(K_i\) around \(m_i\) by
\begin{equation}\label{eq:Ri-Ki-def}
R_i
\;:=\;
\max\Bigl\{\,r_{\mathrm c},\ \frac{2\sigma^2(d-1)}{\alpha_i}\Bigr\},
\qquad
K_i \;:=\; B(m_i,R_i),
\qquad
i\in\{1,2\}.
\end{equation}
We denote by \(p_{\infty,K}^{(i)}\) the conditional measure of \(p_\infty^{(i)}\)
on \(K_i\), that is,
\[
p_{\infty,K}^{(i)}(A)
\;:=\;
\frac{p_\infty^{(i)}(A\cap K_i)}{p_\infty^{(i)}(K_i)},
\qquad
A\subset\mathbb R^d\ \text{Borel}.
\]

Since \(V\) is continuous and \(K_i\) is a compact ball, we have
\[
V_{\min}^{(i)} \;:=\; \inf_{x\in K_i} V(x) \;>\; -\infty,
\qquad
V_{\max}^{(i)} \;:=\; \sup_{x\in K_i} V(x) \;<\; +\infty.
\]
Recall that the invariant density of the diffusion
\[
\mathrm dX_t \;=\; -\nabla V(X_t)\,\mathrm dt + \sigma\,\mathrm dW_t
\]
is (up to normalization) given by
\(
p_\infty(x)\propto \exp\!\bigl(-\tfrac{2}{\sigma^2}V(x)\bigr)
\).
Hence, on \(K_i\),
\begin{equation}\label{eq:density-ratio-Ki}
\frac{\sup_{x\in K_i} p_\infty^{(i)}(x)}{\inf_{x\in K_i} p_\infty^{(i)}(x)}
\;\le\;
\exp\!\Bigl(\frac{2}{\sigma^2}\bigl(V_{\max}^{(i)}-V_{\min}^{(i)}\bigr)\Bigr).
\end{equation}

It is classical that the Lebesgue measure restricted to the ball \(K_i\) satisfies
a Poincar\'e inequality with constant of order \(R_i^2\).
Combining this with the density bound \eqref{eq:density-ratio-Ki}, one obtains
a local Poincar\'e inequality for \(p_{\infty,K}^{(i)}\) on \(K_i\) with explicit
constant; see for instance the discussion before Theorem~1.4 in
Bakry--Barthe--Cattiaux--Guillin~\cite{bakry2008simple}, where the authors
derive the (non-optimal but explicit) bound
\(
\kappa_R \,\lesssim\, R^2 \exp(\mathrm{Osc}_R V)
\)
for \(\mu\) restricted to a ball \(B(0,R)\).
Specialized to our potential \(\tilde V(x)=\tfrac{2}{\sigma^2}V(x)\), this yields
the following lemma.

\begin{lemma}[Local Poincar\'e inequality on \(K_i\)]\label{lem:local-Poincare-Ki}
For each \(i\in\{1,2\}\) and every smooth \(g:\mathbb R^d\to\mathbb R\),
\begin{equation}\label{eq:Ki-Poincare}
\mathrm{Var}_{p_{\infty,K}^{(i)}}(g)
\;\le\;
C_{K_i}\int_{K_i}|\nabla g(x)|^2\,p_{\infty,K}^{(i)}(\mathrm dx),
\end{equation}
where the local Poincar\'e constant can be chosen explicitly as
\begin{equation}\label{eq:CKi-explicit}
C_{K_i}
\;\le\;
c_d\,R_i^2\,
\exp\!\Bigl(\frac{2}{\sigma^2}\bigl(V_{\max}^{(i)}-V_{\min}^{(i)}\bigr)\Bigr),
\end{equation}
for some finite constant \(c_d>0\) depending only on the dimension \(d\),
and \(R_i\) is given in~\eqref{eq:Ri-Ki-def}.
\end{lemma}

\paragraph{Lyapunov function on each well.}
Fix \(i\in\{1,2\}\).
In order to apply the Lyapunov criterion of
Bakry--Barthe--Cattiaux--Guillin~\cite{bakry2008simple}, we construct a
radial Lyapunov function around \(m_i\).
Choose
\begin{equation}\label{eq:gamma-i-def}
\gamma_i
\;:=\;
\frac{\alpha_i}{\sigma^2},
\end{equation}
and define
\begin{equation}\label{eq:Wi-def}
W_i(x)
\;:=\;
\begin{cases}
\exp\bigl(\gamma_i r_i(x)\bigr), & r_i(x)\ge R_i,\vspace{2pt}\\
\text{any $C^2$ extension with }W_i(x)\ge 1, & r_i(x)<R_i,
\end{cases}
\end{equation}
so that \(W_i\in C^2(\Omega_i)\) and \(W_i\ge 1\) everywhere.

On the region \(\{x\in\Omega_i:\ r_i(x)\ge R_i\}\) we have
\[
\nabla W_i(x)
=
\gamma_i\,\frac{x-m_i}{r_i(x)}\,W_i(x),
\qquad
\Delta W_i(x)
=
\Bigl(\gamma_i^2 + \gamma_i\frac{d-1}{r_i(x)}\Bigr)W_i(x).
\]
Recall that the diffusion generator is
\[
L_{\mathrm{diff}} f(x)
\;:=\;
\frac{\sigma^2}{2}\Delta f(x) - \nabla V(x)\cdot\nabla f(x).
\]
Using Assumption~\ref{ass:drift-outside-core}, we compute, for \(r_i(x)\ge R_i\),
\begin{align}
L_{\mathrm{diff}} W_i(x)
&=
\frac{\sigma^2}{2}\Delta W_i(x) - \nabla V(x)\cdot\nabla W_i(x)\nonumber\\
&=
\Biggl[
\frac{\sigma^2}{2}\Bigl(\gamma_i^2 + \gamma_i\frac{d-1}{r_i(x)}\Bigr)
- \gamma_i\,\frac{\langle x-m_i,\nabla V(x)\rangle}{r_i(x)}
\Biggr] W_i(x)\nonumber\\
&\le
\Biggl[
\frac{\sigma^2}{2}\Bigl(\gamma_i^2 + \gamma_i\frac{d-1}{r_i(x)}\Bigr)
- \gamma_i\alpha_i
\Biggr] W_i(x)\nonumber\\
&=
-\gamma_i
\Biggl[
\alpha_i
- \frac{\sigma^2}{2}\gamma_i
- \frac{\sigma^2}{2}\frac{d-1}{r_i(x)}
\Biggr] W_i(x).\label{eq:LWi-pretheta}
\end{align}
With the choice \(\gamma_i=\alpha_i/\sigma^2\), we have
\begin{equation}\label{eq:delta-i-def}
\delta_i
\;:=\;
\alpha_i - \frac{\sigma^2}{2}\gamma_i
=
\frac{\alpha_i}{2}
\;>\;0.
\end{equation}
By the definition of \(R_i\) in~\eqref{eq:Ri-Ki-def},
for all \(x\) with \(r_i(x)\ge R_i\),
\[
\delta_i - \frac{\sigma^2}{2}\frac{d-1}{r_i(x)}
\;\ge\;
\frac{\delta_i}{2}
=
\frac{\alpha_i}{4},
\]
and therefore, combining this with \eqref{eq:LWi-pretheta},
\begin{equation}\label{eq:Lyap-outside-Ki}
L_{\mathrm{diff}} W_i(x)
\;\le\;
-\theta_i\,W_i(x),
\qquad
x\in\Omega_i,\ r_i(x)\ge R_i,
\end{equation}
where we set
\begin{equation}\label{eq:theta-i-def}
\theta_i
\;:=\;
\frac{\gamma_i\delta_i}{2}
=
\frac{1}{2}\cdot\frac{\alpha_i}{\sigma^2}\cdot\frac{\alpha_i}{2}
=
\frac{\alpha_i^2}{4\sigma^2}
\;>\;0.
\end{equation}

On the compact ball \(K_i=B(m_i,R_i)\) we define
\begin{equation}\label{eq:bi-def}
b_i
\;:=\;
\sup_{x\in K_i}
\Bigl(
L_{\mathrm{diff}} W_i(x) + \theta_i W_i(x)
\Bigr)_+,
\qquad
(a)_+ := \max\{a,0\}.
\end{equation}
Since \(K_i\) is compact and \(W_i\in C^2(\Omega_i)\), this supremum is finite.
Combining \eqref{eq:Lyap-outside-Ki} with \eqref{eq:bi-def} yields the global
Lyapunov condition
\begin{equation}\label{eq:Lyap-full-Ki}
L_{\mathrm{diff}} W_i(x)
\;\le\;
-\theta_i\,W_i(x)
\;+\;
b_i\,\mathbf 1_{K_i}(x),
\qquad
\forall x\in\Omega_i.
\end{equation}

\paragraph{Application of the BCG Lyapunov–Poincar\'e criterion.}
We now recall the following result, which is a localized version of
Theorem~1.4 in Bakry--Barthe--Cattiaux--Guillin~\cite{bakry2008simple}.

\begin{theorem}[Lyapunov criterion for Poincar\'e]\label{thm:BCG-Poincare}
Let \(\mu\) be a probability measure on \(\mathbb R^d\) which is reversible
for a diffusion generator \(L\) with carré du champ \(\Gamma\).
Assume there exist a measurable function \(W\ge 1\), constants
\(\theta>0\), \(b<\infty\), and a measurable set \(K\subset\mathbb R^d\) such that
\[
LW \;\le\; -\theta W + b\,\mathbf 1_K
\quad\text{on }\mathbb R^d,
\]
and that the restriction \(\mu_K\) of \(\mu\) to \(K\) satisfies a local
Poincar\'e inequality with constant \(C_K\), i.e.
\[
\mathrm{Var}_{\mu_K}(g)
\;\le\;
C_K \int_K \Gamma(g,g)\,\mathrm d\mu_K,
\qquad
\forall g\in C_c^\infty(\mathbb R^d).
\]
Then \(\mu\) satisfies a global Poincar\'e inequality
\[
\mathrm{Var}_\mu(f)
\;\le\;
\frac{1}{\theta}\bigl(1 + b\,C_K\bigr)
\int \Gamma(f,f)\,\mathrm d\mu,
\qquad
\forall f\in C_c^\infty(\mathbb R^d).
\]
\end{theorem}

We apply Theorem~\ref{thm:BCG-Poincare} to each conditional invariant
measure \(p_\infty^{(i)}\) with the choices
\[
\mu = p_\infty^{(i)},\qquad
L = L_{\mathrm{diff}},\qquad
\Gamma(f,f) = \frac{\sigma^2}{2}|\nabla f|^2,
\]
\[
W = W_i,\qquad
\theta = \theta_i,\qquad
b = b_i,\qquad
K = K_i.
\]
The Lyapunov condition \eqref{eq:Lyap-full-Ki} provides the drift inequality,
while Lemma~\ref{lem:local-Poincare-Ki} provides the local Poincar\'e
inequality on \(K_i\) with constant \(C_{K_i}\).

Therefore, for each \(i\in\{1,2\}\) and every smooth
\(f:\mathbb R^d\to\mathbb R\),
\[
\mathrm{Var}_{p_\infty^{(i)}}(f)
\;\le\;
\frac{1}{\theta_i}\bigl(1 + b_i\,C_{K_i}\bigr)
\int_{\Omega_i} \frac{\sigma^2}{2}|\nabla f(x)|^2\,p_\infty^{(i)}(\mathrm dx).
\]
Comparing with \eqref{eq:local-Poincare-well}, and recalling that the
Dirichlet form is
\(
\mathcal E_{\mathrm{diff}}^{(i)}(f,f)
=
\tfrac{\sigma^2}{2}\int_{\Omega_i}|\nabla f|^2\,p_\infty^{(i)}(\mathrm dx)
\),
we obtain
\begin{equation}\label{eq:lambda-in-i-lower-bound-general}
\lambda_{\mathrm{in}}^{(i)}
\;\ge\;
\frac{\sigma^2\,\theta_i}{2\bigl(1 + b_i\,C_{K_i}\bigr)}.
\end{equation}
Using the explicit value \(\theta_i=\alpha_i^2/(4\sigma^2)\) from
\eqref{eq:theta-i-def}, this becomes
\begin{equation}\label{eq:lambda-in-i-lower-bound-CK}
\lambda_{\mathrm{in}}^{(i)}
\;\ge\;
\frac{\alpha_i^2}{
8\bigl(1 + b_i\,C_{K_i}\bigr)}.
\end{equation}
Finally, plugging the bound \eqref{eq:CKi-explicit} on \(C_{K_i}\), we obtain
the fully explicit estimate
\begin{equation}\label{eq:lambda-in-i-lower-bound-fully-explicit}
\lambda_{\mathrm{in}}^{(i)}
\;\ge\;
\frac{\alpha_i^2}{
8\Bigl(1 + b_i\,c_d\,R_i^2
\exp\!\bigl(\tfrac{2}{\sigma^2}(V_{\max}^{(i)}-V_{\min}^{(i)})\bigr)\Bigr)},
\qquad
R_i = \max\Bigl\{\,r_{\mathrm c},\ \tfrac{2\sigma^2(d-1)}{\alpha_i}\Bigr\}.
\end{equation}

Since
\[
\lambda_{\mathrm{in}}
\;:=\;
\min\bigl\{\lambda_{\mathrm{in}}^{(1)},\lambda_{\mathrm{in}}^{(2)}\bigr\},
\]
we conclude that
\begin{equation}\label{eq:lambda-in-global-explicit}
\lambda_{\mathrm{in}}
\;\ge\;
\min_{i\in\{1,2\}}
\frac{\alpha_i^2}{
8\Bigl(1 + b_i\,c_d\,R_i^2
\exp\!\bigl(\tfrac{2}{\sigma^2}(V_{\max}^{(i)}-V_{\min}^{(i)})\bigr)\Bigr)},
\qquad
R_i = \max\Bigl\{\,r_{\mathrm c},\ \tfrac{2\sigma^2(d-1)}{\alpha_i}\Bigr\}.
\end{equation}
This gives an explicit lower bound on the within-well mixing rate in terms
of the drift constants \(\alpha_i\), the core radius \(r_{\mathrm c}\), the enlarged
radii \(R_i\), the local oscillations \(V_{\max}^{(i)}-V_{\min}^{(i)}\), the noise
level \(\sigma\), and the dimension-dependent constant \(c_d\), without any
convexity assumption on \(V\).


\subsubsection{Controlling Inter-Well Variance: Barrier-Free Coupling via Global Jumps}
\label{sec:lambda-inter}

In the variance decomposition
\begin{equation}\label{eq:Var-decomp}
\mathrm{Var}_{p_\infty}(f)
=
\sum_{i=1}^2 w_i\,\mathrm{Var}_{p_\infty^{(i)}}(f)
+
w_1w_2\bigl(\bar f_{\Omega_1}-\bar f_{\Omega_2}\bigr)^2,
\end{equation}
the first term describes fluctuations \emph{within} each well, while the second term
\[
w_1w_2\bigl(\bar f_{\Omega_1}-\bar f_{\Omega_2}\bigr)^2
\]
measures the difference between the mean values in the two wells.
For a pure diffusion, reducing this term requires transferring mass across the saddle,
so classical bounds involve the barrier height.
In this section we show that, in the presence of a global jump component,
one can \emph{directly jump from one core to the other}, and obtain an
inter-well constant \(\lambda_{\mathrm{inter}}\) which does \emph{not} depend on the saddle height.

Our goal is an inequality of the form
\begin{equation}\label{eq:between-goal}
w_1w_2\bigl(\bar f_{\Omega_1}-\bar f_{\Omega_2}\bigr)^2
\;\le\;
\frac{C_{\mathrm{in}}}{\lambda_{\mathrm{in}}}\,\mathcal E_{\mathrm{diff}}(f,f)
\;+\;
\frac{1}{\lambda_{\mathrm{inter}}}\,\mathcal E_{\mathrm{jump}}(f,f),
\end{equation}
where the first term on the right-hand side depends only on within-well diffusion
(no jumps are needed), and the second term depends only on the global jump part.
We will derive an explicit lower bound on \(\lambda_{\mathrm{inter}}\) in terms of
\(\lambda_0\), the shell density lower bound \(g_\star\), the geometric parameters
\(D,r_{\mathrm c}\), and some local core constants; in particular, no saddle
information enters \(\lambda_{\mathrm{inter}}\).

\paragraph{Global jump structure and shell density lower bound.}
Recall our assumption that the jump measure contains a global component
\[
\nu_{\mathrm{glob}}(\mathrm dr) \;=\; \lambda_0\,\nu_J(\mathrm dr),
\qquad \lambda_0>0,
\]
where \(\nu_J\) is a probability measure on \(\mathbb R^d\) with isotropic density
\[
\nu_J(\mathrm dr)=j(r)\,\mathrm dr,
\qquad
j(r)=j_0(|r|).
\]
Let \(R\sim\nu_J\) and \(S:=|R|\).
The radial density \(g\) of \(S\) is given by
\[
g(s)=|\mathbb S^{d-1}|\,s^{d-1}\,j_0(s),
\qquad s>0,
\]
where \(|\mathbb S^{d-1}|=d\,\omega_d\) and \(\omega_d:=\mathrm{Vol}\bigl(B(0,1)\bigr)\).

Let \(D:=|m_2-m_1|\) be the distance between the two minima and \(r_{\mathrm c}>0\) the core radius.
We assume a uniform lower bound on the radial density on a shell around \(D\):
there exists \(g_\star>0\) such that
\begin{equation}\label{eq:shell-lower}
g(s)\;\ge\;g_\star
\qquad
\forall s\in\bigl[D-2r_{\mathrm c},\,D+2r_{\mathrm c}\bigr].
\end{equation}
Equivalently, since \(g(s)=|\mathbb S^{d-1}|\,s^{d-1}\,j_0(s)\), for all
\(s\in[D-2r_{\mathrm c},\,D+2r_{\mathrm c}]\) we have
\begin{equation}\label{eq:jstar}
j_0(s)
\;\ge\;
\frac{g_\star}{|\mathbb S^{d-1}|(D+2r_{\mathrm c})^{d-1}}
\;=\;
\frac{g_\star}{d\,\omega_d\,(D+2r_{\mathrm c})^{d-1}}
\;=:\; j_\star.
\end{equation}

Geometrically, for any \(x\in\mathcal C_1\), \(y\in\mathcal C_2\),
\[
|y-x|\in\bigl[D-2r_{\mathrm c},\,D+2r_{\mathrm c}\bigr],
\]
hence by \eqref{eq:jstar} we have \(j_0(|y-x|)\ge j_\star\).

\paragraph{A basic Jensen inequality: controlling core mean differences by a cross-core integral.}

Recall that \(\Omega_i\) is the basin of the \(i\)-th well,
\(\mathcal C_i\subset\Omega_i\) is the corresponding core, and
\[
w_i := p_\infty(\Omega_i),\qquad
c_i := p_\infty(\mathcal C_i),
\]
\[
p_\infty^{(i)} := \frac{1}{w_i}p_\infty(\,\cdot\,\cap\Omega_i),
\qquad
p_{\infty,\mathcal C}^{(i)} := \frac{1}{c_i}p_\infty(\,\cdot\,\cap\mathcal C_i).
\]
We write
\[
\bar f_{\Omega_i} := \int_{\Omega_i} f\,\mathrm d p_\infty^{(i)},
\qquad
\bar f_{\mathcal C_i} := \int_{\mathcal C_i} f\,\mathrm d p_{\infty,\mathcal C}^{(i)}.
\]

Let
\[
X\sim p_{\infty,\mathcal C}^{(1)},\qquad
Y\sim p_{\infty,\mathcal C}^{(2)}
\]
be independent.
Then
\[
\mathbb E[f(X)]=\bar f_{\mathcal C_1},\qquad
\mathbb E[f(Y)]=\bar f_{\mathcal C_2}.
\]
Using \(\mathrm{Var}(f(Y)-f(X))\ge 0\), or equivalently Jensen’s inequality
for the convex function \(z\mapsto z^2\), we obtain
\[
\mathbb E\bigl[(f(Y)-f(X))^2\bigr]
\;\ge\;
\bigl(\mathbb E[f(Y)]-\mathbb E[f(X)]\bigr)^2
=
\bigl(\bar f_{\mathcal C_2}-\bar f_{\mathcal C_1}\bigr)^2.
\]
Writing the expectation as an integral, this reads
\begin{equation}\label{eq:jensen-core}
\iint_{\mathcal C_1\times\mathcal C_2}
\bigl(f(y)-f(x)\bigr)^2\,
p_{\infty,\mathcal C}^{(1)}(\mathrm dx)\,
p_{\infty,\mathcal C}^{(2)}(\mathrm dy)
\;\ge\;
\bigl(\bar f_{\mathcal C_2}-\bar f_{\mathcal C_1}\bigr)^2.
\end{equation}
Using
\[
p_{\infty,\mathcal C}^{(1)}(\mathrm dx)
=
\frac{1}{c_1}p_\infty(x)\,\mathbf 1_{\mathcal C_1}(x)\,\mathrm dx,
\qquad
p_{\infty,\mathcal C}^{(2)}(\mathrm dy)
=
\frac{1}{c_2}p_\infty(y)\,\mathbf 1_{\mathcal C_2}(y)\,\mathrm dy,
\]
we can rewrite \eqref{eq:jensen-core} as
\begin{equation}\label{eq:jensen-core-weighted}
\iint_{\mathcal C_1\times\mathcal C_2}
\bigl(f(y)-f(x)\bigr)^2\,
p_\infty(x)p_\infty(y)\,\mathrm dx\,\mathrm dy
\;\ge\;
c_1c_2\bigl(\bar f_{\mathcal C_2}-\bar f_{\mathcal C_1}\bigr)^2.
\end{equation}

This purely measure-theoretic inequality does not use the jump structure
or the potential landscape.
It shows that if we manage to control the cross-core integral on the left
by the jump Dirichlet form, then we can control the difference between the
core means.

\paragraph{Embedding the cross-core integral into the jump Dirichlet form via the shell lower bound.}

Recall the jump Dirichlet form
\[
\mathcal E_{\mathrm{jump}}(f,f)
=
\frac12\int_{\mathbb R^d} p_\infty(x)\int_{\mathbb R^d}
\bigl(f(x+r)-f(x)\bigr)^2\,\nu(\mathrm dr)\,\mathrm dx,
\]
where \(\nu\) contains the global component
\(\nu_{\mathrm{glob}}(\mathrm dr)=\lambda_0\nu_J(\mathrm dr)\).
Let \(\mathcal E_{\mathrm{glob}}\) denote the contribution of \(\nu_{\mathrm{glob}}\),
so that
\[
\mathcal E_{\mathrm{jump}}(f,f)
\;\ge\;
\mathcal E_{\mathrm{glob}}(f,f).
\]

By the change of variables \(y=x+r\) and using the density of \(\nu_J\), we get
\[
\mathcal E_{\mathrm{glob}}(f,f)
=
\frac12\int_{\mathbb R^d} p_\infty(x)\int_{\mathbb R^d}
\bigl(f(y)-f(x)\bigr)^2\,
\lambda_0 j_0(|y-x|)\,\mathrm dy\,\mathrm dx.
\]
Keeping only the contribution from jumps that start in \(\mathcal C_1\) and land in
\(\mathcal C_2\) (discarding other nonnegative terms lowers the right-hand side
and hence yields a lower bound), we obtain
\begin{equation}\label{eq:glob-lower-crosscore}
\mathcal E_{\mathrm{glob}}(f,f)
\;\ge\;
\frac12\int_{\mathcal C_1} p_\infty(x)\int_{\mathcal C_2}
\bigl(f(y)-f(x)\bigr)^2\,
\lambda_0 j_0(|y-x|)\,\mathrm dy\,\mathrm dx.
\end{equation}
For \(x\in\mathcal C_1\), \(y\in\mathcal C_2\) we know
\(|y-x|\in[D-2r_{\mathrm c},D+2r_{\mathrm c}]\), hence by \eqref{eq:jstar},
\(j_0(|y-x|)\ge j_\star\), and therefore
\[
\mathcal E_{\mathrm{glob}}(f,f)
\;\ge\;
\frac12\lambda_0 j_\star
\int_{\mathcal C_1} p_\infty(x)\int_{\mathcal C_2}
\bigl(f(y)-f(x)\bigr)^2\,\mathrm dy\,\mathrm dx.
\]

We now convert the inner \(\mathrm dy\) into a weighted measure \(p_\infty(y)\,\mathrm dy\)
so that we can compare with \eqref{eq:jensen-core-weighted}.
Introduce the local core supremum
\[
M_{\mathcal C}
\;:=\;
\max\Bigl\{\sup_{x\in\mathcal C_1}p_\infty(x),\ \sup_{y\in\mathcal C_2}p_\infty(y)\Bigr\}.
\]
Since \(p_\infty\) is continuous and \(\mathcal C_i\) are compact, \(M_{\mathcal C}<\infty\).
For all \(y\in\mathcal C_2\), \(p_\infty(y)\le M_{\mathcal C}\), hence
\[
\mathrm dy\;\ge\;\frac{1}{M_{\mathcal C}}\,p_\infty(y)\,\mathrm dy.
\]
Plugging this into the previous bound yields
\begin{align}
\mathcal E_{\mathrm{glob}}(f,f)
&\ge
\frac12\lambda_0 j_\star
\int_{\mathcal C_1} p_\infty(x)\int_{\mathcal C_2}
\bigl(f(y)-f(x)\bigr)^2\,\frac{1}{M_{\mathcal C}}p_\infty(y)\,\mathrm dy\,\mathrm dx
\nonumber\\
&=
\frac12\,\lambda_0 j_\star\,\frac{1}{M_{\mathcal C}}
\iint_{\mathcal C_1\times\mathcal C_2}
\bigl(f(y)-f(x)\bigr)^2\,
p_\infty(x)p_\infty(y)\,\mathrm dx\,\mathrm dy.
\label{eq:glob-vs-product}
\end{align}
Combining \eqref{eq:glob-vs-product} with \eqref{eq:jensen-core-weighted}, we obtain
\begin{equation}\label{eq:core-mean-controlled-by-jump}
c_1c_2\bigl(\bar f_{\mathcal C_2}-\bar f_{\mathcal C_1}\bigr)^2
\;\le\;
\frac{2M_{\mathcal C}}{\lambda_0\,j_\star}\,
\mathcal E_{\mathrm{glob}}(f,f)
\;\le\;
\frac{2M_{\mathcal C}}{\lambda_0\,j_\star}\,
\mathcal E_{\mathrm{jump}}(f,f).
\end{equation}

\paragraph{From core means back to well means (purely within-well; no jumps needed).}

The previous step shows that the core mean difference
\(\bar f_{\mathcal C_1}-\bar f_{\mathcal C_2}\) is controlled by the jump energy.
To connect this to the well mean difference
\(\bar f_{\Omega_1}-\bar f_{\Omega_2}\) appearing in \eqref{eq:Var-decomp},
observe that
\[
\bar f_{\Omega_1}-\bar f_{\Omega_2}
=
\bigl(\bar f_{\Omega_1}-\bar f_{\mathcal C_1}\bigr)
+
\bigl(\bar f_{\mathcal C_1}-\bar f_{\mathcal C_2}\bigr)
+
\bigl(\bar f_{\mathcal C_2}-\bar f_{\Omega_2}\bigr),
\]
so by \((a+b+c)^2\le 3(a^2+b^2+c^2)\),
\begin{equation}\label{eq:mean-decomp-three}
\bigl(\bar f_{\Omega_1}-\bar f_{\Omega_2}\bigr)^2
\;\le\;
3\Bigl[
\bigl(\bar f_{\Omega_1}-\bar f_{\mathcal C_1}\bigr)^2
+
\bigl(\bar f_{\mathcal C_1}-\bar f_{\mathcal C_2}\bigr)^2
+
\bigl(\bar f_{\mathcal C_2}-\bar f_{\Omega_2}\bigr)^2
\Bigr].
\end{equation}

We next estimate \(\bar f_{\Omega_i}-\bar f_{\mathcal C_i}\) using within-well variance.
Let
\(
\eta_i := p_\infty^{(i)}(\mathcal C_i)=c_i/w_i
\)
be the mass of the core in the \(i\)-th well.
A simple estimate (e.g.\ by Cauchy–Schwarz or a one-point Poincaré inequality)
gives
\begin{equation}\label{eq:Omega-core-mean-bound}
\bigl(\bar f_{\Omega_i}-\bar f_{\mathcal C_i}\bigr)^2
\;\le\;
\frac{1}{\eta_i}\,
\mathrm{Var}_{p_\infty^{(i)}}(f),
\qquad
i=1,2.
\end{equation}
Insert \eqref{eq:Omega-core-mean-bound} into \eqref{eq:mean-decomp-three},
multiply both sides by \(w_1w_2\), and recall that
\(w_1w_2=(c_1c_2)/(\eta_1\eta_2)\).
Using \eqref{eq:core-mean-controlled-by-jump}, we obtain
\begin{align}
w_1w_2\bigl(\bar f_{\Omega_1}-\bar f_{\Omega_2}\bigr)^2
&\le
3w_1w_2\Bigl[
\frac{1}{\eta_1}\mathrm{Var}_{p_\infty^{(1)}}(f)
+
\frac{1}{\eta_2}\mathrm{Var}_{p_\infty^{(2)}}(f)
\Bigr]
+
3w_1w_2\bigl(\bar f_{\mathcal C_1}-\bar f_{\mathcal C_2}\bigr)^2
\nonumber\\
&\le
3\Bigl(\frac{w_1w_2}{\eta_1}+\frac{w_1w_2}{\eta_2}\Bigr)
\sum_{i=1}^2\mathrm{Var}_{p_\infty^{(i)}}(f)
+
\frac{6M_{\mathcal C}}{\lambda_0j_\star}\,
\mathcal E_{\mathrm{jump}}(f,f).\label{eq:between-pre-lambda}
\end{align}
Define
\begin{equation}\label{eq:Cin-def}
C_{\mathrm{in}}
:=3\Bigl(\frac{w_1w_2}{\eta_1}+\frac{w_1w_2}{\eta_2}\Bigr).
\end{equation}

On the other hand, from the within-well analysis we already have a lower bound
on the within-well mixing rate \(\lambda_{\mathrm{in}}\) such that
\begin{equation}\label{eq:Step-in}
\sum_{i=1}^2 w_i\,\mathrm{Var}_{p_\infty^{(i)}}(f)
\;\le\;
\frac{1}{\lambda_{\mathrm{in}}}\,\mathcal E_{\mathrm{diff}}(f,f).
\end{equation}
Since all weights \(w_i,\eta_i\) are fixed positive constants determined by the
equilibrium measure and the choice of cores, there exists a finite constant
\(C_{\mathrm{in}}>0\) (which can be chosen as in \eqref{eq:Cin-def} up to a
universal factor) such that
\[
\sum_{i=1}^2\mathrm{Var}_{p_\infty^{(i)}}(f)
\;\le\;
\frac{C_{\mathrm{in}}}{\lambda_{\mathrm{in}}}\,\mathcal E_{\mathrm{diff}}(f,f).
\]
Plugging this into \eqref{eq:between-pre-lambda} and enlarging
\(C_{\mathrm{in}}\) if necessary, we obtain
\begin{equation}\label{eq:between-controlled}
w_1w_2\bigl(\bar f_{\Omega_1}-\bar f_{\Omega_2}\bigr)^2
\;\le\;
\frac{C_{\mathrm{in}}}{\lambda_{\mathrm{in}}}\,
\mathcal E_{\mathrm{diff}}(f,f)
\;+\;
\frac{1}{\lambda_{\mathrm{inter}}}\,
\mathcal E_{\mathrm{jump}}(f,f),
\end{equation}
where
\begin{equation}\label{eq:lambda-inter-explicit}
\lambda_{\mathrm{inter}}
\;:=\;
\frac{\lambda_0\,j_\star}{6M_{\mathcal C}}\;\eta_1\eta_2.
\end{equation}

Structurally, \(\lambda_{\mathrm{inter}}\) depends only on
\(\lambda_0\), \(j_\star\) (hence on \(g_\star,D,r_{\mathrm c}\)),
and the local quantities \(M_{\mathcal C},\eta_1,\eta_2\);
\emph{it does not involve the saddle position or the barrier height.}

\paragraph{ Explicit expression for \texorpdfstring{$\lambda_{\mathrm{inter}}$}{lambda\_inter} in terms of the shell lower bound \texorpdfstring{$g_\star$}{g\_*}.}

From \eqref{eq:jstar} we have
\[
j_\star
=
\frac{g_\star}{d\,\omega_d\,(D+2r_{\mathrm c})^{d-1}}.
\]
Substituting this into \eqref{eq:lambda-inter-explicit} yields
\begin{equation}\label{eq:lambda-inter-shell}
\lambda_{\mathrm{inter}}
\;=\;
\frac{\lambda_0}{6M_{\mathcal C}}\;\eta_1\eta_2\,
\frac{g_\star}{d\,\omega_d\,(D+2r_{\mathrm c})^{d-1}}.
\end{equation}
Thus, small values of \(\lambda_{\mathrm{inter}}\) can arise from:
\begin{itemize}
  \item The geometric “angular penalty” factor
  \(
  (D+2r_{\mathrm c})^{-(d-1)}
  \),
  reflecting the small solid angle occupied by the target core in high dimensions;
  \item A small shell density lower bound \(g_\star\), i.e.\ if the jump
  distribution places little radial mass near radius \(D\).
\end{itemize}
Crucially, there is no exponential suppression in terms of the potential barrier
height; the bottleneck for inter-well motion is entirely controlled by jump
geometry and shell density.

\subsubsection{Combining Within-Well and Inter-Well Bounds: A Two-Scale Poincaré Inequality}
\label{sec:two-scale-PI}

We now combine the within-well bound \eqref{eq:Step-in} with the inter-well
bound \eqref{eq:between-controlled} in the variance decomposition
\eqref{eq:Var-decomp}.
We obtain
\begin{align*}
\mathrm{Var}_{p_\infty}(f)
&=
\sum_{i=1}^2 w_i\,\mathrm{Var}_{p_\infty^{(i)}}(f)
+
w_1w_2\bigl(\bar f_{\Omega_1}-\bar f_{\Omega_2}\bigr)^2
\\
&\le
\frac{1}{\lambda_{\mathrm{in}}}\,\mathcal E_{\mathrm{diff}}(f,f)
+
\frac{C_{\mathrm{in}}}{\lambda_{\mathrm{in}}}\,
\mathcal E_{\mathrm{diff}}(f,f)
+
\frac{1}{\lambda_{\mathrm{inter}}}\,
\mathcal E_{\mathrm{jump}}(f,f)
\\
&=
\frac{1+C_{\mathrm{in}}}{\lambda_{\mathrm{in}}}\,\mathcal E_{\mathrm{diff}}(f,f)
+
\frac{1}{\lambda_{\mathrm{inter}}}\,
\mathcal E_{\mathrm{jump}}(f,f).
\end{align*}
If we redefine
\(
\tilde\lambda_{\mathrm{in}}:=\lambda_{\mathrm{in}}/(1+C_{\mathrm{in}})
\)
and, for simplicity of notation, relabel \(\tilde\lambda_{\mathrm{in}}\) again
as \(\lambda_{\mathrm{in}}\) (this only changes the constant by a fixed factor),
we may rewrite the inequality as
\[
\mathrm{Var}_{p_\infty}(f)
\;\le\;
\frac{1}{\lambda_{\mathrm{in}}}\,\mathcal E_{\mathrm{diff}}(f,f)
+
\frac{1}{\lambda_{\mathrm{inter}}}\,\mathcal E_{\mathrm{jump}}(f,f).
\]
Since
\(
\mathcal E_{\mathrm{diff}}(f,f)\le\mathcal E_{\mathrm{levy}}(f,f)
\)
and
\(
\mathcal E_{\mathrm{jump}}(f,f)\le\mathcal E_{\mathrm{levy}}(f,f)
\)
(because \(\mathcal E_{\mathrm{levy}}=\mathcal E_{\mathrm{diff}}+\mathcal E_{\mathrm{jump}}\)),
we arrive at a two–scale Poincaré inequality
\begin{equation}\label{eq:two-scale-PI}
\mathrm{Var}_{p_\infty}(f)
\;\le\;
\Bigl(\frac{1}{\lambda_{\mathrm{in}}}+\frac{1}{\lambda_{\mathrm{inter}}}\Bigr)
\mathcal E_{\mathrm{levy}}(f,f).
\end{equation}
This implies the following lower bound on the spectral gap \(\lambda(\mathcal G)\)
of the generator \(\mathcal G\):
\begin{equation}\label{eq:harmonic}
\lambda(\mathcal G)
\;\ge\;
\Bigl(\lambda_{\mathrm{in}}^{-1}+\lambda_{\mathrm{inter}}^{-1}\Bigr)^{-1}.
\end{equation}

\subsubsection{Harmonic Mean Form: The Overall Rate is Governed by the Slower Scale (up to a Factor \texorpdfstring{$\tfrac12$}{1/2})}

For any \(a,b>0\),
\begin{equation}\label{eq:harmonic-vs-min}
\bigl(a^{-1}+b^{-1}\bigr)^{-1}
=
\frac{1}{\frac{1}{a}+\frac{1}{b}}
=
\frac{ab}{a+b}.
\end{equation}
Since
\(
a+b\le 2\max\{a,b\}
\),
we deduce
\[
\frac{ab}{a+b}
\;\ge\;
\frac{ab}{2\max\{a,b\}}
=
\frac12\min\{a,b\},
\]
that is,
\begin{equation}\label{eq:harmonic-min}
\bigl(a^{-1}+b^{-1}\bigr)^{-1}
\;\ge\;
\frac12\min\{a,b\}.
\end{equation}
Applying \eqref{eq:harmonic-min} with
\(a=\lambda_{\mathrm{in}},\,b=\lambda_{\mathrm{inter}}\)
and combining with \eqref{eq:harmonic}, we obtain the more “intuitive” bound
\begin{equation}\label{eq:min-form}
\lambda(\mathcal G)
\;\ge\;
\Bigl(\lambda_{\mathrm{in}}^{-1}+\lambda_{\mathrm{inter}}^{-1}\Bigr)^{-1}
\;\ge\;
\frac12\min\{\lambda_{\mathrm{in}},\lambda_{\mathrm{inter}}\}.
\end{equation}
This shows that, up to a factor \(1/2\), the global convergence rate is governed
by the slower of the two mechanisms:
\begin{itemize}
  \item If within-well mixing is slow (\(\lambda_{\mathrm{in}}\) small), then even
  very strong jumps cannot improve the overall spectral gap beyond this scale;
  \item If within-well mixing is fast but inter-well jumps are too weak
  (\(\lambda_{\mathrm{inter}}\) small), then mass cannot be exchanged efficiently
  between wells, and the overall convergence remains slow.
\end{itemize}

\subsubsection{Relation to the Saddle: What “Barrier-Free” Means Here}

Along this line of argument we have derived two key constants:
\begin{itemize}
  \item \(\lambda_{\mathrm{in}}\), the within-well spectral gap, controlled purely
  by local information near the minima (e.g.\ Hessian lower bounds or
  Lyapunov conditions);
  \item \(\lambda_{\mathrm{inter}}\), the inter-well rate generated by global jumps,
  depending only on the jump intensity \(\lambda_0\), the radial shell density
  lower bound \(g_\star\), the geometry \((D,r_{\mathrm c})\), and the local
  core constants \(M_{\mathcal C},\eta_1,\eta_2\).
\end{itemize}
At no point did we use the saddle location or the barrier height.
In particular, the expression \eqref{eq:lambda-inter-shell} contains no factor
of the form \(\exp(-H/\sigma^2)\).
Instead, the inter-well bottleneck is entirely controlled by the probability
that a global jump lands in the other core, encoded by shell geometry and
the density lower bound \(g_\star\).

In words:
\begin{quote}
For a pure diffusion, inter-well moves must follow the energy landscape and
typically go through the saddle, so the slowest step involves climbing the
barrier and the spectral gap is often exponentially small in the barrier height.

With a suitably chosen global jump component, inter-well transitions no longer
have to respect the potential landscape: they are governed by the jump
distribution at radius \(\approx D\), i.e.\ by the shell density and geometry.
The metastable “barrier” is effectively bypassed, and the inter-well rate
\(\lambda_{\mathrm{inter}}\) is \emph{barrier-free} in the sense that it is
independent of the saddle height.
\end{quote}

\subsubsection{Two-Scale Poincar\'e Inequality and Harmonic-Mean Lower Bound}\label{sec:two-scale-final}

Combining the exact variance decomposition \eqref{eq:var-decomp-Omega},
the within-well bound \eqref{eq:within-controlled-Omega},
and the inter-well bound encoded in \eqref{eq:lambda-inter-explicit},
we obtain a two-scale Poincar\'e inequality of the schematic form
\[
\mathrm{Var}_{p_\infty}(f)
\;\le\;
\Bigl(\frac{1}{\lambda_{\mathrm{in}}}+\frac{1}{\lambda_{\mathrm{inter}}}\Bigr)\,
\mathcal E_{\mathrm{levy}}(f,f),
\]
and hence the spectral gap satisfies
\begin{equation}\label{eq:hmean-gap}
\lambda(\mathcal G)
\;\ge\;
\Bigl(\lambda_{\mathrm{in}}^{-1}+\lambda_{\mathrm{inter}}^{-1}\Bigr)^{-1}.
\end{equation}

Finally, using the elementary inequality
\[
(a^{-1}+b^{-1})^{-1}
=
\frac{ab}{a+b}
\;\ge\;
\frac12\min\{a,b\},
\]
we conclude the intuitive bound
\begin{equation}\label{eq:min-gap}
\lambda(\mathcal G)
\;\ge\;
\frac12\min\{\lambda_{\mathrm{in}},\lambda_{\mathrm{inter}}\}.
\end{equation}


